\begin{definition}[Complete zipper fusion family]%
    \label{def:mpdo_complete_zipper_fusion_family}
    \lean{MPOTensor.CompleteZipperFusionFamily}
    \leanok
    \uses{def:mpo_tensor}
    Let $\Lambda$ be a finite set. A \emph{complete zipper fusion family}
    consists of positive bond dimensions $D_a$, MPO blocks
    $B_a^{ij}\in\MN{D_a}$, fusion multiplicities $N_{ab}^{c}$,
    and matrices
    \begin{align}
        X^c_{ab,\mu}    & :\C^{D_c}\longrightarrow
        \C^{D_a}\otimes\C^{D_b}, \notag            \\
        X^{c+}_{ab,\mu} & :\C^{D_a}\otimes\C^{D_b}
        \longrightarrow\C^{D_c}. \notag
    \end{align}
    The matrices $X^{c+}_{ab,\mu}$ are independently chosen left inverses,
    not adjoints. For all labels and multiplicity indices,
    \begin{align}
        X^{d+}_{ab,\nu}X^c_{ab,\mu}
         & =\delta_{cd}\delta_{\mu\nu}\Id_{D_c}.
        \label{eq:rfp_fusion_biorthogonality}
    \end{align}
    The opposite product
    $P_{ab}=\sum_{c,\mu}X^c_{ab,\mu}X^{c+}_{ab,\mu}$ is the projector onto
    the product-MPO support; it need not be the identity on the ambient bond
    space or an orthogonal projector. Write
    $B_{ab}^{ik}=\sum_j B_a^{ij}\otimes B_b^{jk}$.
    Every product-MPO matrix is reconstructed on this support:
    \begin{align}
        B_{ab}^{ik}
         & =\sum_{c,\mu}X^c_{ab,\mu}B_c^{ik}X^{c+}_{ab,\mu}.
        \label{eq:rfp_fusion_reconstruction}
    \end{align}
    The two zipper identities are
    \begin{align}
        B_{ab}^{ik}X^c_{ab,\mu}
         & =X^c_{ab,\mu}B_c^{ik},
        \label{eq:rfp_zipper_synthesis} \\
        X^{c+}_{ab,\mu}B_{ab}^{ik}
         & =B_c^{ik}X^{c+}_{ab,\mu}.
        \label{eq:rfp_zipper_analysis}
    \end{align}
    Each block is injective, and the matrix of coefficients of all labelled
    blocks has a simultaneous left inverse. Contraction of its row indexed
    by $(c,x,y)$ with $B_d^{ij}$ gives $E_{xy}$ when $d=c$ and zero otherwise.
    These are the fusion-tensor, reconstruction, zipper, and simultaneous
    inverse relations of
    \cite[lines~161, 181--200, and 269--277]{Bultinck2017Anyons}.
\end{definition}

\begin{theorem}[Complete zipper family from a BNT]%
    \label{thm:mpdo_blocked_bnt_complete_zipper}
    \lean{MPOTensor.BNTFusionIsometryFamily.blockedFusionSynthesis}
    \lean{MPOTensor.BNTFusionIsometryFamily.blockedFusionAnalysis}
    \lean{MPOTensor.BNTFusionIsometryFamily.%
        exists_blocked_completeZipperFusionFamily}
    \leanok
    \uses{thm:mpdo_bnt_common_positive_blocked_fusion, thm:mpdo_pair_fusion_bnt_completeness, thm:mpdo_one_letter_simultaneous_inverse, def:mpdo_complete_zipper_fusion_family}
    Suppose that the labelled tensors of a fusion-isometry family with
    length-independent structure coefficients
    $c_{\alpha,\beta,\gamma}^{(L)}
        =\tr(\chi_{\alpha,\beta,\gamma}^{L})$
    form a BNT. There is a common positive blocking length at which they
    determine a complete zipper fusion family. Its
    bond dimensions are the original labelled bond dimensions, its fusion
    multiplicities are
    $N_{\alpha\beta}^{\gamma}=\dim\chi_{\alpha,\beta,\gamma}$,
    its synthesis map is the adjoint of $U_{\alpha,\beta}$, and its analysis
    map is $U_{\alpha,\beta}$, after the canonical identification
    $\C^{D_\alpha D_\beta}
        \cong\C^{D_\alpha}\otimes\C^{D_\beta}$.
    The two zipper identities, exact reconstruction, injectivity, and the
    simultaneous left inverse of the block coefficients follow as conclusions
    rather than additional hypotheses.
\end{theorem}
\begin{proof}\leanok
    Length independence makes each positive diagonal matrix
    $\chi_{\alpha,\beta,\gamma}$ the identity. Use a common positive
    blocking length for the BNT and carry the fusion equation through
    the blocking. Completeness of the pair decomposition gives
    $U_{\alpha,\beta}U_{\alpha,\beta}^{\dagger}=\Id$, so the synthesis
    map followed by the analysis map is the identity. The other three
    matrix identities are the blocked zipper identities
    (\ref{eq:rfp_zipper_synthesis}) and
    (\ref{eq:rfp_zipper_analysis}) and the reconstruction identity
    (\ref{eq:rfp_fusion_reconstruction}). At this blocking length, the
    labelled matrices span the direct sum of the block matrix algebras.
    This supplies both block injectivity and the simultaneous left inverse
    of the block coefficients. These maps satisfy
    Definition~\ref{def:mpdo_complete_zipper_fusion_family}.
\end{proof}

\begin{theorem}[Complete zipper family from split compressions]%
    \label{thm:mpdo_complete_zipper_of_compression}
    \lean{MPOTensor.CompleteZipperFusionFamily.ofCompression}
    \lean{MPOTensor.CompleteZipperFusionFamily.ofCompression_fusionTensor}
    \lean{MPOTensor.CompleteZipperFusionFamily.%
        ofCompression_fusionTensorLeftInverse}
    \lean{MPOTensor.CompleteZipperFusionFamily.%
        pairLetter_eq_compressionSynthesis_mul}
    \lean{MPOTensor.CompleteZipperFusionFamily.%
        compressionAnalysis_mul_compressionSynthesis}
    \lean{MPOTensor.CompleteZipperFusionFamily.%
        exists_pairCompression_remainder_eq_zero_of_star}
    \lean{MPOTensor.CompleteZipperFusionFamily.ofStar}
    \lean{MPOTensor.CompleteZipperFusionFamily.%
        exists_ofStar_fusionTensor_eq}
    \leanok
    \uses{def:mpdo_complete_zipper_fusion_family, thm:asym_multiblock}
    Let $B_a^{ij}$, $a\in\Lambda$, be injective MPO blocks of positive bond
    dimension whose labelled letters have a simultaneous left inverse as in
    Definition~\ref{def:mpdo_complete_zipper_fusion_family}. Suppose that for
    every pair $a,b$ the product tensor $B_{ab}$ admits a multi-block
    compression, as in Theorem~\ref{thm:asym_multiblock}, onto $N_{ab}^{c}$
    copies of each block $B_c$, with vanishing remainder. Then the
    compression maps $V_{(c,\mu)}$ and $W_{(c,\mu)}$ are the fusion tensors
    $X^c_{ab,\mu}$ and their left inverses $X^{c+}_{ab,\mu}$ of a complete
    zipper fusion family. The vanishing remainder is the absence of nonzero
    blocks above the diagonal assumed in
    \cite[lines~181--191]{Bultinck2017Anyons}.
    In particular, if for every nonempty word the word trace of every product
    $B_{ab}$ is the sum of the word traces of $N_{ab}^{c}$ copies of each
    $B_c$, and $(B_{ab}^{ik})^\dagger$ lies in the unital algebra generated by
    the $B_{ab}^{jl}$ for every $i,k$, then such compressions exist and the
    conclusion holds.
\end{theorem}
\begin{proof}\leanok
    \uses{def:mpdo_complete_zipper_fusion_family, thm:asym_multiblock,
        thm:asym_splitting}
    Biorthogonality (\ref{eq:rfp_fusion_biorthogonality}) is
    $W_sV_t=\delta_{st}\Id$. Since the remainder
    $B_{ab}^{ik}-\sum_sV_sB_{c(s)}^{ik}W_s$ vanishes, this sum is the
    reconstruction (\ref{eq:rfp_fusion_reconstruction}). Multiplying the
    reconstruction by $X^c_{ab,\mu}$ on the right, or by $X^{c+}_{ab,\mu}$ on
    the left, and using biorthogonality gives the two zipper identities.
    In the star-closed case the word-algebra module of each product is semisimple, so
    its multi-block compression can be chosen with vanishing remainder.
\end{proof}

\begin{theorem}[Simultaneous left inverse after blocking]%
    \label{thm:mpdo_complete_zipper_of_compression_blocked}
    \lean{MPSTensor.exists_positive_wordTupleSpanTop_of_isNormal}
    \lean{MPSTensor.wordTupleSpanTop_of_smul}
    \lean{MPOTensor.exists_blockTensor_isMPOBlockLeftInverse}
    \lean{MPSTensor.MultiBlockCompression.ofWords}
    \lean{MPSTensor.MultiBlockCompression.remainder_ofWords_eq_zero}
    \lean{MPOTensor.CompleteZipperFusionFamily.PairCompression.blockTensor}
    \lean{MPOTensor.CompleteZipperFusionFamily.jointBlockLength}
    \lean{MPOTensor.CompleteZipperFusionFamily.ofCompressionBlocked}
    \lean{MPOTensor.CompleteZipperFusionFamily.ofCompressionBlocked_tensor}
    \lean{MPOTensor.CompleteZipperFusionFamily.%
        ofCompressionBlocked_fusionTensor}
    \lean{MPOTensor.CompleteZipperFusionFamily.%
        ofCompressionBlocked_fusionTensorLeftInverse}
    \lean{MPOTensor.CompleteZipperFusionFamily.ofStarBlocked}
    \leanok
    \uses{thm:mpdo_complete_zipper_of_compression, def:mpdo_block_left_inverse,
        def:blockwise_product_word_span}
    Let $B_1,\ldots,B_m$ be MPO blocks of positive bond dimension which,
    read as tensors on the doubled physical alphabet, are normal, and no two
    of which are gauge equivalent up to a nonzero scalar. Then there is $L>0$
    such that the tuples $(B_1^w,\ldots,B_m^w)$ of the words $w$ of length
    $L$ span $\bigoplus_cM_{D_c}$, as in
    Definition~\ref{def:blockwise_product_word_span}, so the blocked tensors
    $B_c^{[L]}$ are injective and have a simultaneous left inverse at one blocked site, as in
    Definition~\ref{def:mpdo_block_left_inverse}: there are coefficients
    $C_{(c,x,y),(I,J)}$ with
    $\sum_{I,J}C_{(c,x,y),(I,J)}\bigl((B_d^{[L]})^{IJ}\bigr)_{x'y'}
        =\delta_{cd}\delta_{xx'}\delta_{yy'}$. If
    every product $B_aB_b$ admits a multi-block compression onto $N_{ab}^c$
    copies of each $B_c$ with vanishing remainder, as in
    Theorem~\ref{thm:mpdo_complete_zipper_of_compression}, then the same
    compression maps are the fusion tensors of a complete zipper fusion
    family of the blocked tensors $B_c^{[L]}$, with no simultaneous left
    inverse among the hypotheses. Under the same hypotheses on the blocks,
    this applies in the star-closed case of
    Theorem~\ref{thm:mpdo_complete_zipper_of_compression}.
    The simultaneous left inverse at one unblocked site, as printed
    in~\cite[line~269]{Bultinck2017Anyons}, can fail for injective,
    inequivalent blocks, and the scalar in the non-equivalence cannot be
    dropped~\cite{gap:bmwshv17_joint_block_left_inverse}.
\end{theorem}
\begin{proof}\leanok
    \uses{thm:mpdo_complete_zipper_of_compression, def:mpdo_block_left_inverse,
        def:blockwise_product_word_span, thm:mpdo_one_letter_simultaneous_inverse}
    After a gauge and a nonzero rescaling each block has transfer map of
    spectral radius one. Rescaling a block multiplies its words of length $L$
    by a nonzero constant, so it changes neither the joint span nor
    non-equivalence up to a scalar. For normal, pairwise inequivalent
    tensors the block-injectivity argument
    of~\cite[lines~317--345]{Cirac2016MPDO_arXiv} gives one length $L>0$ at
    which the word tuples span $\bigoplus_cM_{D_c}$. Grouping $L$ sites into
    one physical index turns the words of length $L$ into single letters.
    Projecting the span onto the summand $c$ shows that the letters of
    $B_c^{[L]}$ span $M_{D_c}$, so each blocked tensor is injective.
    The coefficients $C_{(c,x,y),(I,J)}$ expressing the matrix unit
    $E_{xy}$ of the summand $c$, with zero in every other summand, as a
    combination of the blocked letters are the left inverse of
    Definition~\ref{def:mpdo_block_left_inverse}, by
    Theorem~\ref{thm:mpdo_one_letter_simultaneous_inverse}. Each blocked
    letter of a product is a word of the unblocked letters. In the
    coordinates $g$ of a split compression every unblocked letter is block
    diagonal, $g\,B_{ab}^{ik}g^{-1}=\bigoplus_s B_{c(s)}^{ik}$, with one
    summand for each slot $s=(c,\mu)$. Hence for the word $w=(i_1k_1)\cdots
        (i_Lk_L)$ of the blocked letter $(I,J)$,
    \begin{align}
        g\,(B_{ab}^{[L]})^{IJ}g^{-1}
         & =\prod_{t=1}^{L}\Bigl(\bigoplus_sB_{c(s)}^{i_tk_t}\Bigr)
        =\bigoplus_s\prod_{t=1}^{L}B_{c(s)}^{i_tk_t}
        =\bigoplus_s(B_{c(s)}^{[L]})^{IJ}.
        \notag
    \end{align}
    The compression
    therefore remains split after blocking, with the same compression maps,
    and Theorem~\ref{thm:mpdo_complete_zipper_of_compression} applies to the
    blocked tensors.
\end{proof}

\begin{definition}[Final-block selection from one physical index pair]
    \label{def:mpdo_complete_zipper_final_block_selector}
    \lean{MPOTensor.CompleteZipperFusionFamily.finalBlockSelector}
    \leanok
    \uses{def:mpdo_complete_zipper_fusion_family}
    Fix a final label $d$. Summing the diagonal matrix-unit coefficients of
    the simultaneous left inverse defines
    \begin{align}
        s_d(i,j)
         & =\sum_{x=0}^{D_d-1}C_{(d,x,x),(i,j)}. \notag
    \end{align}
    Here $C_{(d,x,y),(i,j)}$ denotes the coefficient of the simultaneous
    left inverse corresponding to the matrix unit $E_{xy}$ in block $d$.
\end{definition}

\begin{lemma}[Final-block selection identity]
    \label{lem:mpdo_complete_zipper_final_block_selector}
    \lean{MPOTensor.CompleteZipperFusionFamily.finalBlockSelector_apply,
        MPOTensor.CompleteZipperFusionFamily.finalBlockSelector_sum}
    \leanok
    \uses{def:mpdo_complete_zipper_final_block_selector}
    The coefficients $s_d$ satisfy
    \begin{align}
        \sum_{i,j}s_d(i,j)B_e^{ij}
         & =
        \begin{cases}
            \Id_{D_d}, & e=d,    \\
            0,         & e\ne d.
        \end{cases} \notag
    \end{align}
    Thus contraction over one physical index pair selects the identity on
    the chosen final block and annihilates every other final block.
\end{lemma}
\begin{proof}\leanok
    Apply the simultaneous left inverse entrywise and sum over the diagonal
    matrix units of block $d$. If $e=d$, their sum is $\Id_{D_d}$; if
    $e\ne d$, every term vanishes.
\end{proof}

\begin{theorem}[F-move and inverse]%
    \label{thm:mpdo_complete_zipper_printed_fmove}
    \lean{MPOTensor.CompleteZipperFusionFamily.%
        rightTripleSynthesis_mul_printedFMatrix}
    \lean{MPOTensor.CompleteZipperFusionFamily.%
        printedFMatrix_mul_inversePrintedFMatrix}
    \lean{MPOTensor.CompleteZipperFusionFamily.%
        inversePrintedFMatrix_mul_printedFMatrix}
    \leanok
    \uses{def:mpdo_complete_zipper_fusion_family, lem:mpdo_complete_zipper_final_block_selector, lem:matrix_kronecker_identity_cancellation}
    Fix $a,b,c,d\in\Lambda$. There is an invertible matrix
    \begin{align}
        F^{abc}_d:
        \bigoplus_e\C^{N_{ab}^{e}}\otimes\C^{N_{ec}^{d}}
         & \longrightarrow
        \bigoplus_f\C^{N_{bc}^{f}}\otimes\C^{N_{af}^{d}} \notag
    \end{align}
    whose rows are indexed by $(f,\lambda,\sigma)$ and whose columns are
    indexed by $(e,\mu,\nu)$. Its entries satisfy
    \begin{align}
        (X^e_{ab,\mu}\otimes\Id)X^d_{ec,\nu}
         & =
        \sum_{f,\lambda,\sigma}
        (F^{abc}_d)^{f\lambda\sigma}_{e\mu\nu}
        (\Id\otimes X^f_{bc,\lambda})X^d_{af,\sigma}.
        \label{eq:rfp_printed_fmove}
    \end{align}
    The matrix in the opposite orientation is its two-sided inverse. This is
    the $F$-move of \cite[Section~3.4]{Bultinck2017Anyons}.
\end{theorem}
\begin{proof}\leanok
    \uses{def:mpdo_complete_zipper_final_block_selector, lem:matrix_kronecker_identity_cancellation}
    Form the full left- and right-associated synthesis matrices $L$ and $R$
    and their analysis matrices $L^+$ and $R^+$. Applying
    (\ref{eq:rfp_fusion_reconstruction}) twice to the two associative
    decompositions gives
    \begin{align}
        B_{abc}^{il}
         & =L D_{\mathrm L}^{il}L^+
        =R D_{\mathrm R}^{il}R^+.
        \label{eq:rfp_triple_reconstruction}
    \end{align}
    Contract (\ref{eq:rfp_triple_reconstruction}) with the coefficients
    $s_d$ from
    Definition~\ref{def:mpdo_complete_zipper_final_block_selector} and sum
    over the final labels $d$. The direct sums $D_{\mathrm L}$ and
    $D_{\mathrm R}$ both contract to the identity. Consequently,
    \begin{align}
        P    & :=LL^+=RR^+,
        \label{eq:rfp_triple_support} \\
        L^+L & =R^+R=\Id. \notag
    \end{align}
    Hence $R^+L$ changes left-tree coordinates into right-tree coordinates and
    \begin{align}
        R(R^+L)
         & =(RR^+)L=PL=(LL^+)L=L. \notag
    \end{align}
    Its opposite-oriented comparison is $L^+R$. Using
    (\ref{eq:rfp_triple_support}),
    \begin{align}
        (R^+L)(L^+R) & =R^+PR=\Id, \notag \\
        (L^+R)(R^+L) & =L^+PL=\Id. \notag
    \end{align}
    The zipper identities (\ref{eq:rfp_zipper_synthesis}) and
    (\ref{eq:rfp_zipper_analysis}) show that $R^+L$ intertwines the direct
    sums of the final-block matrices. The simultaneous left inverse
    annihilates its blocks with different final labels. On the
    $d$-diagonal block, injectivity gives
    $(R^+L)_d=F^{abc}_d\otimes\Id_{D_d}$.
    Restricting $R(R^+L)=L$ to this block gives
    (\ref{eq:rfp_printed_fmove}).
    The corresponding fixed-final block of $L^+R$, together with
    $(R^+L)(L^+R)=\Id=(L^+R)(R^+L)$, gives the inverse identities
    tensored with $\Id_{D_d}$.
    Lemma~\ref{lem:matrix_kronecker_identity_cancellation} removes the
    nonempty bond-space identity factor and yields the two inverse identities
    for $F^{abc}_d$.
\end{proof}

\begin{figure}[ht]
    \centering
    % Each internal-edge label records the intermediate charge together with
    % its multiplicity coordinate; the root-edge label records the final
    % charge and the multiplicity coordinate of the second fusion.
    \begin{align}
        \tntree{((a\,b)_{e,\mu}\,c)_{d,\nu}}
         & =
        \sum_{f,\lambda,\sigma}
        (F^{abc}_{d})^{f\lambda\sigma}_{e\mu\nu}
        \tntree{(a\,(b\,c)_{f,\lambda})_{d,\sigma}}. \notag
    \end{align}
    \caption{The source-oriented $F$-move. The left-associated fusion tensor
        is expanded in the right-associated fusion basis. The upper indices
        $(f,\lambda,\sigma)$ label the rows of $F^{abc}_d$, and the lower
        indices $(e,\mu,\nu)$ label its columns. This is the printed-$F$
        row/column convention used here for
        \cite[Section~3.4]{Bultinck2017Anyons}.}
    \label{fig:mpdo_complete_zipper_printed_fmove}
\end{figure}

\begin{theorem}[Gauge covariance of the F-matrices]%
    \label{thm:mpdo_complete_zipper_fmatrix_gauge}
    \lean{MPOTensor.CompleteZipperFusionFamily.regauge}
    \lean{MPOTensor.CompleteZipperFusionFamily.regauge_fusionTensor}
    \lean{MPOTensor.CompleteZipperFusionFamily.%
        eq_printedFMatrix_of_rightTripleSynthesis_mul}
    \lean{MPOTensor.CompleteZipperFusionFamily.%
        rightTreeGauge_mul_printedFMatrix_regauge}
    \lean{MPOTensor.CompleteZipperFusionFamily.printedFMatrix_regauge}
    \lean{MPOTensor.CompleteZipperFusionFamily.inversePrintedFMatrix_regauge}
    \leanok
    \uses{def:mpdo_complete_zipper_fusion_family, thm:mpdo_complete_zipper_printed_fmove,
        lem:matrix_kronecker_identity_cancellation}
    For every triple of labels let $Y^c_{ab}$ be an invertible
    $N_{ab}^{c}\times N_{ab}^{c}$ matrix, write
    $\Pi_{ab}=\bigoplus_c Y^c_{ab}\otimes\Id_{D_c}$, and let $X_{ab}$ and
    $X^+_{ab}$ be the synthesis and analysis matrices whose blocks are the
    $X^c_{ab,\mu}$ and the $X^{c+}_{ab,\mu}$. Replace them by
    $X_{ab}\Pi_{ab}$ and $\Pi_{ab}^{-1}X^+_{ab}$. The new fusion tensors are
    \begin{align}
        X'^c_{ab,\mu}
         & =\sum_\nu (Y^c_{ab})_{\nu\mu}X^c_{ab,\nu}. \notag
    \end{align}
    The new matrices, with the same blocks, multiplicities and simultaneous
    left inverse, again form a complete zipper fusion family. The matrix
    $F^{abc}_d$ is the only matrix satisfying (\ref{eq:rfp_printed_fmove}),
    and the matrix $F'^{abc}_d$ of the new family is
    \begin{align}
        F'^{abc}_d
         & =G_{\mathrm R}^{-1}F^{abc}_d G_{\mathrm L},
         &
        G_{\mathrm L}
         & =\bigoplus_e Y^e_{ab}\otimes Y^d_{ec},
         &
        G_{\mathrm R}
         & =\bigoplus_f Y^f_{bc}\otimes Y^d_{af}.
        \label{eq:rfp_fmatrix_gauge}
    \end{align}
    Its inverse transforms as
    $(F'^{abc}_d)^{-1}=G_{\mathrm L}^{-1}(F^{abc}_d)^{-1}G_{\mathrm R}$.
    This is the gauge freedom of the fusion tensors in
    \cite[lines~164--166]{Bultinck2017Anyons} and the induced transformation of
    the $F$-symbols in \cite[lines~415--422]{GarreRubio2023MPOSym}.
\end{theorem}
\begin{proof}\leanok
    \uses{def:mpdo_complete_zipper_fusion_family, thm:mpdo_complete_zipper_printed_fmove,
        lem:matrix_kronecker_identity_cancellation}
    Write $\mathcal B^{ik}_{ab}=\bigoplus_c\Id_{N_{ab}^c}\otimes B_c^{ik}$,
    so that (\ref{eq:rfp_fusion_reconstruction}) reads
    $B_{ab}^{ik}=X_{ab}\mathcal B^{ik}_{ab}X^+_{ab}$. Since $\Pi_{ab}$ acts
    only on the multiplicity factors, it commutes with
    $\mathcal B^{ik}_{ab}$. The biorthogonality
    (\ref{eq:rfp_fusion_biorthogonality}) and the reconstruction
    (\ref{eq:rfp_fusion_reconstruction}) of the new family are
    \begin{align}
        (\Pi_{ab}^{-1}X^+_{ab})(X_{ab}\Pi_{ab})
         & =\Pi_{ab}^{-1}(X^+_{ab}X_{ab})\Pi_{ab}=\Id,
        \notag                                                      \\
        (X_{ab}\Pi_{ab})\mathcal B^{ik}_{ab}(\Pi_{ab}^{-1}X^+_{ab})
         & =X_{ab}\mathcal B^{ik}_{ab}\Pi_{ab}\Pi_{ab}^{-1}X^+_{ab}
        =X_{ab}\mathcal B^{ik}_{ab}X^+_{ab}=B_{ab}^{ik}.
        \notag
    \end{align}
    The zipper identities (\ref{eq:rfp_zipper_synthesis}) and
    (\ref{eq:rfp_zipper_analysis}) of the new family are
    \begin{align}
        B_{ab}^{ik}X_{ab}\Pi_{ab}
         & =X_{ab}\mathcal B^{ik}_{ab}\Pi_{ab}
        =(X_{ab}\Pi_{ab})\mathcal B^{ik}_{ab},
        \notag                                        \\
        \Pi_{ab}^{-1}X^+_{ab}B_{ab}^{ik}
         & =\Pi_{ab}^{-1}\mathcal B^{ik}_{ab}X^+_{ab}
        =\mathcal B^{ik}_{ab}(\Pi_{ab}^{-1}X^+_{ab}).
        \notag
    \end{align}
    If $R(G\otimes\Id)=L$ for the fixed-final synthesis matrices $L$ and $R$,
    multiplying on the left by $R^+$ and using $R^+R=\Id$ gives
    $G\otimes\Id=F^{abc}_d\otimes\Id$, hence $G=F^{abc}_d$.
    Expanding the composites shows that the new synthesis matrices are
    $L'=L(G_{\mathrm L}\otimes\Id)$ and $R'=R(G_{\mathrm R}\otimes\Id)$.
    The $F$-move of the new family becomes
    $R\bigl(G_{\mathrm R}F'^{abc}_d\otimes\Id\bigr)=L(G_{\mathrm L}\otimes\Id)
        =R\bigl(F^{abc}_dG_{\mathrm L}\otimes\Id\bigr)$.
    Cancelling $R$ gives $G_{\mathrm R}F'^{abc}_d=F^{abc}_dG_{\mathrm L}$,
    which is (\ref{eq:rfp_fmatrix_gauge}). The inverse formula follows because
    two-sided inverses are unique.
\end{proof}

\begin{definition}[Fourfold fusion maps and changes of basis]%
    \label{def:mpdo_complete_zipper_fourfold_fusion_data}
    \lean{MPOTensor.CompleteZipperFusionFamily.%
        FourfoldLeftAssocMultiplicity}
    \lean{MPOTensor.CompleteZipperFusionFamily.%
        FourfoldLeftInnerMultiplicity}
    \lean{MPOTensor.CompleteZipperFusionFamily.FourfoldMiddleMultiplicity}
    \lean{MPOTensor.CompleteZipperFusionFamily.FourfoldPairMultiplicity}
    \lean{MPOTensor.CompleteZipperFusionFamily.FourfoldRightAssocMultiplicity}
    \lean{MPOTensor.CompleteZipperFusionFamily.leftAssocFourfoldSynthesis}
    \lean{MPOTensor.CompleteZipperFusionFamily.leftInnerFourfoldSynthesis}
    \lean{MPOTensor.CompleteZipperFusionFamily.middleFourfoldSynthesis}
    \lean{MPOTensor.CompleteZipperFusionFamily.pairFourfoldSynthesis}
    \lean{MPOTensor.CompleteZipperFusionFamily.rightAssocFourfoldSynthesis}
    \lean{MPOTensor.CompleteZipperFusionFamily.%
        leftAssocToLeftInnerPrintedFMatrix}
    \lean{MPOTensor.CompleteZipperFusionFamily.%
        leftInnerToMiddlePrintedFMatrix}
    \lean{MPOTensor.CompleteZipperFusionFamily.%
        middleToRightAssocPrintedFMatrix}
    \lean{MPOTensor.CompleteZipperFusionFamily.leftAssocToPairPrintedFMatrix}
    \lean{MPOTensor.CompleteZipperFusionFamily.pairToRightAssocPrintedFMatrix}
    \lean{MPOTensor.CompleteZipperFusionFamily.threeEdgePrintedFMatrix}
    \lean{MPOTensor.CompleteZipperFusionFamily.twoEdgePrintedFMatrix}
    \lean{MPOTensor.CompleteZipperFusionFamily.%
        leftInnerToLeftAssocInversePrintedFMatrix}
    \lean{MPOTensor.CompleteZipperFusionFamily.%
        middleToLeftInnerInversePrintedFMatrix}
    \lean{MPOTensor.CompleteZipperFusionFamily.%
        rightAssocToMiddleInversePrintedFMatrix}
    \lean{MPOTensor.CompleteZipperFusionFamily.%
        pairToLeftAssocInversePrintedFMatrix}
    \lean{MPOTensor.CompleteZipperFusionFamily.%
        rightAssocToPairInversePrintedFMatrix}
    \lean{MPOTensor.CompleteZipperFusionFamily.%
        threeEdgeInversePrintedFMatrix}
    \lean{MPOTensor.CompleteZipperFusionFamily.%
        twoEdgeInversePrintedFMatrix}
    \leanok
    \uses{def:mpdo_complete_zipper_fusion_family, thm:mpdo_complete_zipper_printed_fmove}
    Fix labels $a,b,c,d,e$. The five fusion trees in equation~(33) of
    \cite{Bultinck2017Anyons} determine five multiplicity spaces and five
    synthesis maps into the fourfold bond space. These spaces use the integer
    multiplicities $N_{ab}^{c}$ of the complete zipper family, rather than
    multiplicity coordinates weighted by the positive diagonal matrices
    $\chi_{\alpha,\beta,\gamma}$.

    Each edge of the pentagon carries the $F$-matrix for the affected
    triple tensor factor and the identity on the untouched multiplicity
    factor. The three upper edges and the two lower edges define composites
    in the forward orientation. Replacing each edge by its inverse defines
    the opposite-oriented composites whose entries occur in
    equation~(33).
\end{definition}

\begin{figure}[ht]
    \centering
    % M sits on top, LI and R in the middle row, L and P on the bottom.
    % The upper route L -> LI -> M -> R climbs the left flank; the lower
    % route L -> P -> R crosses the bottom, its labels swapped outward.
    % Tree subscripts are precisely the intermediate fusion charges
    % carried by the corresponding internal edges.
    \begin{tikzcd}[column sep=tiny]
        & \tntree{(a\,((b\,c)_h\,d)_i)_e} \arrow[dr, "F^{bcd}_{i}"] & \\
        \tntree{((a\,(b\,c)_h)_g\,d)_e} \arrow[ur, "F^{ahd}_{e}"] & &
        \tntree{(a\,(b\,(c\,d)_j)_i)_e} \\
        \tntree{(((a\,b)_f\,c)_g\,d)_e} \arrow[u, "F^{abc}_{g}"]
        \arrow[rr, "F^{fcd}_{e}"'] & &
        \tntree{((a\,b)_f\,(c\,d)_j)_e} \arrow[u, "F^{abj}_{e}"']
    \end{tikzcd}
    \caption{The five complete zipper fusion trees of
        \cite[Section~3.4, equation~(33)]{Bultinck2017Anyons}. Each arrow is
        a printed $F$-move, with right-tree coordinates indexing rows and
        left-tree coordinates indexing columns.}
    \label{fig:mpdo_complete_zipper_fusion_pentagon}
\end{figure}

\begin{theorem}[Complete zipper fusion pentagon]%
    \label{thm:mpdo_complete_zipper_fusion_pentagon}
    \lean{MPOTensor.CompleteZipperFusionFamily.%
        threeEdgePrintedFMatrix_eq_twoEdgePrintedFMatrix}
    \lean{MPOTensor.CompleteZipperFusionFamily.%
        threeEdgeInversePrintedFMatrix_eq_twoEdgeInversePrintedFMatrix}
    \lean{MPOTensor.CompleteZipperFusionFamily.inversePrintedFMatrix_pentagon}
    \leanok
    \uses{def:mpdo_complete_zipper_fourfold_fusion_data, thm:mpdo_complete_zipper_printed_fmove, lem:matrix_kronecker_identity_cancellation}
    The three-edge and two-edge composites of the forward
    $F$-matrices are equal. The corresponding composites of the inverse
    $F$-matrices are also equal. In the upper/lower convention of
    \cite[equation~(33)]{Bultinck2017Anyons}, the latter equality is
    \begin{align}
         & \sum_{h,\sigma,\lambda,\omega}
        (F^{abc}_{g})^{f\mu\nu}_{h\sigma\lambda}
        (F^{ahd}_{e})^{g\lambda\rho}_{i\omega\kappa}
        (F^{bcd}_{i})^{h\sigma\omega}_{j\gamma\delta} \notag \\
         & \hspace{3em}=
        \sum_{\tau}
        (F^{fcd}_{e})^{g\nu\rho}_{j\gamma\tau}
        (F^{abj}_{e})^{f\mu\tau}_{i\delta\kappa}.
        \label{eq:rfp_fusion_pentagon}
    \end{align}
    Every entry in (\ref{eq:rfp_fusion_pentagon}) is an entry of the inverse
    of the forward matrix in (\ref{eq:rfp_printed_fmove}), not the reciprocal
    of an individual entry.
\end{theorem}
\begin{proof}\leanok
    \uses{lem:matrix_kronecker_identity_cancellation}
    Apply the elementary three-block $F$-move along each of the five edges.
    Both paths give the same expansion in the right-associated fourfold
    fusion map. Its analysis map is a left inverse, by three successive
    applications of fusion-tensor biorthogonality
    (\ref{eq:rfp_fusion_biorthogonality}). This gives equality after
    tensoring with the identity on the nonzero final bond space, and
    Lemma~\ref{lem:matrix_kronecker_identity_cancellation} gives equality of
    the forward composites. The reverse composites are two-sided inverses of
    this common matrix, hence they coincide. Taking a matrix entry gives
    (\ref{eq:rfp_fusion_pentagon}).
\end{proof}

\begin{theorem}[Coherence of four bond-coordinate reassociations]%
    \label{thm:mpdo_fourfold_bond_coordinate_coherence}
    \lean{MPOTensor.BNTFusionIsometryFamily.%
        fourfoldBondReassociation_threeEdge_eq_twoEdge}
    \lean{MPOTensor.BNTFusionIsometryFamily.%
        fourfoldBondMatrix_threeEdge_eq_twoEdge}
    \leanok
    \uses{def:mpdo_operator_product_associator, def:mpdo_bnt_fusion_isometry_family}
    Fix four labels $\alpha,\beta,\gamma,\delta$ in a BNT fusion-isometry
    family, and write $A=\{0,\ldots,D_\alpha-1\}$,
    $B=\{0,\ldots,D_\beta-1\}$, $C=\{0,\ldots,D_\gamma-1\}$, and
    $D=\{0,\ldots,D_\delta-1\}$ for their bond-coordinate sets. Let
    \begin{align}
        r_{X,Y,Z}:((X\times Y)\times Z)
         & \longrightarrow X\times(Y\times Z) \notag
    \end{align}
    denote the canonical reassociation. Define
    \begin{align}
        r_1 & =r_{A,B,C}\times\Id_D, \notag  \\
        r_2 & =r_{A,B\times C,D}, \notag     \\
        r_3 & =\Id_A\times r_{B,C,D}, \notag \\
        s_1 & =r_{A\times B,C,D}, \notag     \\
        s_2 & =r_{A,B,C\times D}. \notag
    \end{align}
    The three-edge and two-edge reassociations from
    $(((A\times B)\times C)\times D)$ to
    $A\times(B\times(C\times D))$ coincide:
    \begin{align}
        r_3\circ r_2\circ r_1
         & =s_2\circ s_1.
        \label{eq:rfp_bond_reassociation}
    \end{align}
    If $P(r)_{x,y}=1$ when $r(x)=y$ and is zero otherwise, then the
    corresponding coordinate-permutation matrices satisfy
    \begin{align}
        P(r_1)P(r_2)P(r_3)
         & =P(s_1)P(s_2).
        \label{eq:rfp_bond_permutation}
    \end{align}
    These are equalities of Cartesian bond-coordinate identifications only;
    they are distinct from the fusion-multiplicity identity in
    Theorem~\ref{thm:mpdo_complete_zipper_fusion_pentagon}.
\end{theorem}
\begin{proof}\leanok
    Every element $(((a,b),c),d)$ is sent by either composite to
    $(a,(b,(c,d)))$. This proves (\ref{eq:rfp_bond_reassociation}). The
    identity $P(v\circ u)=P(u)P(v)$ then turns
    (\ref{eq:rfp_bond_reassociation}) into
    (\ref{eq:rfp_bond_permutation}).
\end{proof}

\begin{figure}[ht]
    \centering
    % M sits on top, LI and R in the middle row, L and P on the bottom;
    % r_1,r_2,r_3 trace the three-edge route L -> LI -> M -> R, while
    % s_1,s_2 trace L -> P -> R.  The vertices are Cartesian
    % parenthesizations, not fusion trees.
    \begin{tikzcd}[column sep=small]
        & (A\times((B\times C)\times D)) \arrow[dr, "r_3"] & \\
        ((A\times(B\times C))\times D) \arrow[ur, "r_2"] & &
        (A\times(B\times(C\times D))) \\
        (((A\times B)\times C)\times D) \arrow[u, "r_1"] \arrow[rr, "s_1"'] & &
        ((A\times B)\times(C\times D)) \arrow[u, "s_2"']
    \end{tikzcd}
    \caption{The five parenthesizations of four bond-coordinate factors.
        The upper path is the three-edge composite
        $r_3\circ r_2\circ r_1$, and the lower path is the two-edge composite
        $s_2\circ s_1$, depicting canonical Cartesian reassociation of the
        bond coordinates.}
    \label{fig:mpdo_fourfold_bond_coordinate_reassociation}
\end{figure}
